% !TEX program = lualatex
\documentclass[11pt,a4paper]{article}
\usepackage{fontspec}
\usepackage[russian]{babel}
\setmainfont{Tinos}
\usepackage[a4paper,margin=25mm]{geometry}
\usepackage{amssymb}
\begin{document}
\thispagestyle{empty}
\begin{flushright}
Иван Проскураков\\
независимый исследователь\\
\texttt{[email placeholder]}\\[1em]
\today
\end{flushright}

\vspace{2em}
\noindent Редакционная коллегия математического журнала

\vspace{1.5em}
\noindent Уважаемые коллеги,

\vspace{0.8em}
\noindent на рассмотрение предлагается рукопись «Фазово-секторные сертификаты и конструктивное декодирование цифровых разложений: универсальный формализм и экспериментальная демонстрация на числе $\pi$».

Работа вводит универсальную фазово-секторную модель для произвольного вещественного носителя $\alpha\in\mathbb R$, в которой цифровая строка рассматривается как сектор единичного интервала, а позиция - как фазовая точка орбиты $\{B^M\alpha\}$. Число $\pi$ используется как основной экспериментальный пример, но формализм не зависит от специальных свойств $\pi$. Основной положительный результат состоит в построении конечного сертифицирующего locate-контура и в формулировке внешнего секторного сертификата, позволяющего алгебраически извлекать следующую цифру или блок цифр при наличии независимого уточнения фазы.  Важное уточнение: метод является конструктивным декодером и proof-object архитектурой, а не статистическим предиктором; он переводит независимую фазовую информацию из одного представления числа в другое. В обновлённой версии отдельно подчёркивается, что такой внешний сертификат не является мистическим источником данных: он может быть вычислен алгебраически из независимых формул для $\pi$, например из формулы Мэчина или интервальных алгоритмов.

К рукописи приложено экспериментальное приложение, отражающее исследовательскую программу 001--018. Оно включает как положительные результаты, так и отрицательные проверки простых резонансных, drift- и spiral-предикторов. Это сделано для честного разграничения подтверждённого конечного контура и открытой задачи автономного предсказания за горизонтом.

Рукопись не утверждает нормальность $\pi$, не доказывает существование блока из $1000$ нулей и не заявляет вычисления новых цифр за публичным мировым горизонтом без построения соответствующего внешнего сертификата. Вместо этого предлагается воспроизводимая вычислительная модель с точными пропорциональными сертификатами и кросс-валидацией различных представлений вещественных чисел.

\vspace{2em}
\noindent С уважением,\\[1em]
Иван Проскураков
\end{document}
